\section{Conclusion}
\label{sec:conclusion}

This paper presented a systematic analysis of noise robustness in encoder-only transformer architectures, identifying critical vulnerability transitions at layers 3 and 8 that correspond to boundaries between linguistic processing phases. Through evaluation of 300,000 perturbed samples across five models, we demonstrated that these transitions represent universal computational properties, with 61.1\% correlation in vulnerability patterns across architectures.

Our key findings include: (1) RoBERTa's superior robustness (98.8\%) stems from training choices aligning with phase boundaries, with real-world noise (OCR, social media) proving 15-20\% more challenging than synthetic perturbations; (2) Strategic layer dropout achieves 2.47× measured speedup (2.8× at batch=32) while maintaining 95\% accuracy, validating theoretical predictions; (3) The 61.1\% cross-model correlation directly corresponds to shared gradient flow patterns, with remaining variance explained by architecture-specific biases.

Empirical validation confirms theoretical predictions—mutual information measurements show inflection points at layers 3 and 8, gradient norms exhibit 2.3× peaks at transitions, and phase boundaries align with linguistic processing hierarchy. Preliminary GPT-2 experiments reveal decoder transitions at layers 4 and 10, shifted due to causal attention constraints.

\textbf{Practical Implications:} For deployment in noise-critical applications, we recommend RoBERTa-based architectures, implementation of quality-aware routing for adaptive processing, and targeted denoising at identified vulnerable layers. These strategies can reduce computational costs while maintaining robustness.

\textbf{Future Directions:} Important research areas include: (1) Systematic analysis of decoder architectures to identify generation-specific vulnerabilities; (2) Multilingual studies to determine universality of transitions; (3) Development of phase-aware architectures that explicitly model transition boundaries; (4) Runtime validation of theoretical efficiency gains in production systems.

The identification of universal phase transitions advances our understanding of transformer architectures beyond black-box models toward interpretable systems with predictable vulnerability patterns. As transformer models become increasingly critical in real-world applications, this knowledge enables development of more robust and efficient NLP systems that can reliably handle the noisy, imperfect data encountered in practice.